\documentclass[11pt,oneside,openany]{book}
\usepackage[T1]{fontenc}
\usepackage[utf8]{inputenc}
\usepackage[sc]{mathpazo}
\renewcommand{\sfdefault}{lmss}
\renewcommand{\ttdefault}{lmtt}
\usepackage{amsmath,amssymb,amsthm,mathtools}
\usepackage[paperwidth=7.5in,paperheight=10in,inner=.85in,outer=.85in,top=.85in,bottom=.85in,headheight=15pt,headsep=20pt]{geometry}
\usepackage{microtype}
\usepackage{xcolor,graphicx,booktabs,array,longtable}
\newcolumntype{L}[1]{>{\raggedright\arraybackslash}p{#1}}
\usepackage{tikz,pgfplots}
\usetikzlibrary{arrows.meta,calc,patterns,decorations.markings}
\pgfplotsset{compat=1.18}
\definecolor{ink}{HTML}{17334B}
\definecolor{teal}{HTML}{137C83}
\definecolor{copper}{HTML}{AE5932}
\definecolor{soft}{HTML}{EEF3F5}
\pgfplotsset{every axis/.append style={axis line style={black!65},tick style={black!60},tick label style={font=\small},label style={font=\small},legend style={font=\small,draw=none,fill=white},grid style={black!10},major grid style={black!10},line width=.6pt},every axis plot/.append style={line width=1.15pt}}
\usepackage{enumitem,needspace,etoolbox}
\setlist[enumerate]{topsep=5pt,itemsep=5pt,parsep=0pt}
\setlist[itemize]{topsep=4pt,itemsep=3pt,parsep=0pt}
\usepackage{titlesec,fancyhdr}
\titleformat{\chapter}[display]{\normalfont\color{ink}}{\large\sffamily\MakeUppercase{Chapter \thechapter}}{12pt}{\huge\bfseries\raggedright\hyphenpenalty=10000\exhyphenpenalty=10000}[\vspace{12pt}\titlerule]
\titlespacing*{\chapter}{0pt}{-10pt}{24pt}
\titleformat{\section}{\large\bfseries\color{ink}}{\thesection}{.65em}{}
\titleformat{\subsection}{\normalsize\bfseries\color{ink}}{\thesubsection}{.6em}{}
\titlespacing*{\section}{0pt}{18pt plus 3pt minus 2pt}{7pt}
\titlespacing*{\subsection}{0pt}{12pt plus 2pt minus 2pt}{5pt}
\pagestyle{fancy}
\fancyhf{}
\fancyhead[L]{\small\nouppercase{\leftmark}}
\fancyfoot[C]{\small\thepage}
\renewcommand{\headrulewidth}{.3pt}
\renewcommand{\chaptermark}[1]{\markboth{\thechapter\quad #1}{}}
\fancypagestyle{plain}{\fancyhf{}\fancyfoot[C]{\small\thepage}\renewcommand{\headrulewidth}{0pt}}
\usepackage[font=small,labelfont=bf,labelsep=period]{caption}
\usepackage{float,placeins}
\floatplacement{figure}{htbp}
\setcounter{topnumber}{3}
\setcounter{bottomnumber}{2}
\renewcommand{\topfraction}{.9}
\renewcommand{\textfraction}{.08}
\renewcommand{\floatpagefraction}{.8}
\usepackage{makeidx}
\usepackage[columns=3,font=small,indentunit=11pt,columnsep=14pt]{idxlayout}
\makeindex
\usepackage[pdfencoding=auto,psdextra,colorlinks=true,linkcolor=ink,citecolor=teal,urlcolor=teal,bookmarksnumbered=true]{hyperref}
\usepackage{bookmark}
\hypersetup{pdftitle={The Calculus of Probability Distributions},pdfsubject={An undergraduate introduction to information geometry}}
\urlstyle{same}
\newtheorem{theorem}{Theorem}[chapter]
\newtheorem{proposition}[theorem]{Proposition}
\newtheorem{lemma}[theorem]{Lemma}
\newtheorem{corollary}[theorem]{Corollary}
\theoremstyle{definition}
\newtheorem{definition}[theorem]{Definition}
\newtheorem{example}{Example}[chapter]
% A worked example may continue naturally across a page break.  Space, rather
% than a heavy colored box, distinguishes it from the surrounding discussion.
\BeforeBeginEnvironment{example}{\par\addvspace{9pt}\needspace{4\baselineskip}}
\AfterEndEnvironment{example}{\par\addvspace{8pt}}
\theoremstyle{remark}
\newtheorem*{remark}{Remark}
\newenvironment{solution}{\par\smallskip\noindent\textit{Solution.}\ }{\par\smallskip}
\newenvironment{exercises}{\begin{enumerate}[label=\thechapter.\arabic*,ref=\thechapter.\arabic*,leftmargin=2.8em,itemsep=6pt]}{\end{enumerate}}
\newenvironment{answers}{\begin{description}[leftmargin=2.9em,labelwidth=2.5em,labelsep=.4em,itemsep=8pt,font=\normalfont\bfseries]}{\end{description}}
\newcommand{\R}{\mathbb{R}}
\newcommand{\E}{\mathbb{E}}
\newcommand{\Var}{\operatorname{Var}}
\newcommand{\Cov}{\operatorname{Cov}}
\newcommand{\KL}[2]{D\!\left(#1\,\middle\|\,#2\right)}
\newcommand{\grad}{\operatorname{grad}}
\newcommand{\arcosh}{\operatorname{arcosh}}
\newcommand{\sech}{\operatorname{sech}}
\newcommand{\softmax}{\operatorname{softmax}}
\newcommand{\diag}{\operatorname{diag}}
\newcommand{\dd}{\,\mathrm{d}}
\newcommand{\ip}[2]{\left\langle #1,#2\right\rangle}
\newcommand{\p}{\mathbf{p}}
\newcommand{\q}{\mathbf{q}}
\newcommand{\f}{\mathbf{f}}
\newcommand{\uvec}{\mathbf{u}}
\newcommand{\vvec}{\mathbf{v}}
\newcommand{\ones}{\mathbf{1}}
\newcommand{\simp}{\Delta^{n-1}}
\newcommand{\simpi}{\Delta^{n-1}_{\!\circ}}
\numberwithin{equation}{chapter}
\setlength{\parindent}{1.2em}
\setlength{\parskip}{0pt}
\linespread{1.115}
\setlength{\emergencystretch}{2em}
\clubpenalty=10000
\widowpenalty=10000
\displaywidowpenalty=10000
\raggedbottom
\setcounter{tocdepth}{1}
\setcounter{secnumdepth}{2}
% Leave enough room for the four-character subsection numbers in Chapter 10.
\makeatletter
\renewcommand*\l@section{\@dottedtocline{1}{1.5em}{3.35em}}
\makeatother


\hypersetup{pdftitle={Chapter 9: Natural Gradients and Potential-Driven Evolution}}
\begin{document}
\mainmatter
\setcounter{chapter}{8}
\chapter{Natural Gradients and Potential-Driven Evolution}
\label{ch:natural}

A scalar objective assigns a value to every statistical state. Its differential specifies how that value changes along a velocity. A metric converts the differential into a gradient vector, selecting the direction of steepest increase for the metric's definition of unit length. With the Fisher metric, the resulting natural gradient has the replicator form.

The familiar coexistence equation from Chapter 1 provides a guide. Its equilibrium attracted nearby proportions, and later chapters described that motion using entropy and Fisher length. We can now ask which scalar function increases as the state moves and how the metric converts that function into the observed velocity. The answer connects the geometry to the original dynamics.

\section{The metric definition of a gradient}\index{gradient flow}
Let $\Phi$ be a differentiable function on a domain with metric $g$. Its differential $d\Phi_p$ is the linear function on velocities that gives directional derivatives. The metric gradient is the unique vector satisfying
\begin{equation}\label{eq:gradientdefinition}
g_p(\grad_g\Phi,v)=d\Phi_p(v)
\quad\text{for every tangent vector }v.
\end{equation}
Uniqueness follows from positive definiteness. In coordinates, if $d\Phi(v)=(\nabla\Phi)^Tv$ and $g(u,v)=u^TGv$, then
\begin{equation}\label{eq:metricgradient}
\grad_g\Phi=G^{-1}\nabla\Phi.
\end{equation}
The coordinate derivatives form a covector, a linear rule for evaluating velocities. The inverse metric turns this rule into a tangent vector.

Cauchy--Schwarz shows that among all vectors with $g(v,v)=1$, the largest value of $d\Phi(v)$ is $\|\grad_g\Phi\|_g$. It is attained in the normalized gradient direction when the gradient is nonzero. Thus the defining identity recovers the geometric meaning of steepest ascent.

For the Fisher metric the gradient is called the \emph{natural gradient}\index{natural gradient}. The associated continuous ascent equation is $p'=\grad_F\Phi$. Optimization of a loss function instead uses the negative natural gradient.

\section{The Bernoulli gradient and the original equation}
In one dimension, the metric is $ds^2=g(x)dx^2$, so \eqref{eq:gradientdefinition} gives gradient component $\Phi'(x)/g(x)$. The Bernoulli metric therefore gives
\begin{equation}\label{eq:binarynatural}
(\grad_F\Phi)^x=x(1-x)\Phi'(x).
\end{equation}
Every smooth scalar advantage $\Delta(x)$ has an antiderivative $\Phi(x)=\int^x\Delta(s)ds$. The corresponding binary replicator equation is its Fisher-gradient ascent.

\begin{example}[A potential for stable coexistence]
Take $\Phi(x)=3x-2x^2$.
\begin{solution}
Its derivative is $3-4x$. Hence natural-gradient ascent is
\[
x'=x(1-x)(3-4x),
\]
the recurring coexistence equation. The potential has its unique interior maximum at $x=3/4$, with value $9/8$. Along the flow,
\[
\dot\Phi=x(1-x)(3-4x)^2\geq0.
\]
The endpoints are equilibria of the smooth extended replicator field; the gradient is defined on the positive interval.
\end{solution}
\end{example}

\begin{figure}[htbp]
\centering
\begin{tikzpicture}
\begin{axis}[width=.86\textwidth,height=5cm,xmin=0,xmax=1,ymin=0,ymax=1.25,xlabel={probability $x$},ylabel={potential $\Phi(x)$},grid=major,xtick={0,.25,.5,.75,1}]
\addplot[teal,domain=0:1,samples=80]{3*x-2*x^2};
\addplot[only marks,ink] coordinates {(.75,1.125)};
\end{axis}
\end{tikzpicture}
\caption{The potential generating the coexistence model under the Bernoulli Fisher metric. Interior trajectories move toward its maximum.}
\end{figure}

The Fisher multiplier changes the rate of ascent at different probabilities. Near an interior equilibrium, it rescales the ordinary derivative field by a positive factor. Near the boundary it becomes small, reflecting the local length cost of moving a rare category's probability.

\section{Deriving the simplex natural gradient}
Let $\Phi(\p)$ be smooth on the positive simplex, and choose a smooth extension to its normalization neighborhood. Write $f_i=\partial_i\Phi$ and $\bar f=\sum_i p_if_i$. Consider
\begin{equation}\label{eq:simplexnatural}
\boxed{(\grad_F\Phi)_i=p_i\left(\partial_i\Phi-\sum_jp_j\partial_j\Phi\right).}
\end{equation}
The components sum to zero. For every tangent vector $v$,
\begin{align*}
g_\p(\grad_F\Phi,v)
&=\sum_i(f_i-\bar f)v_i\\
&=\sum_if_iv_i-\bar f\sum_iv_i=d\Phi_\p(v).
\end{align*}
It therefore satisfies the gradient definition.

The extension does not affect the result. Two extensions agreeing on the simplex have the same directional derivatives on its tangent plane. Their ambient gradients differ there by a multiple of $\ones$, which cancels under the weighted mean subtraction. The resulting vector depends only on the function restricted to the simplex.

\begin{example}[A linear objective]
Let $\Phi=2p_A+p_B$ and $\p=(1/3,1/3,1/3)$. Find the natural gradient and its ascent rate.
\begin{solution}
The ambient derivative vector is $(2,1,0)$, with probability-weighted mean 1. The natural gradient is $(1/3,0,-1/3)$. The directional derivative is $2(1/3)+0=2/3$. Its squared Fisher norm is $3(1/9+1/9)=2/3$.
\end{solution}
\end{example}

For comparison within the same constrained problem, the Euclidean tangent gradient of this objective is $(1,0,-1)$, obtained by subtracting the arithmetic mean of its ambient derivative components. Following that constant vector from the uniform point reaches the boundary in finite time. The Fisher gradient adjusts with the probabilities, and its constant-fitness solution remains positive at every finite time.

\section{Ascent, variance, and symmetric interactions}\index{potential}\index{symmetric payoff matrix}
Natural-gradient ascent gives a particularly direct relation between the objective and the velocity. The rate at which the objective increases is the squared Fisher speed. For replicator equations that speed is already known from Chapter 7: it is fitness variance. Thus one identity can be read as a statement about optimization, motion, or variation among the types.
Along natural-gradient ascent,
\begin{equation}\label{eq:ascentidentity}
\boxed{\dot\Phi=d\Phi(\grad_F\Phi)
=\|\grad_F\Phi\|_F^2=\Var_\p(f)\geq0.}
\end{equation}
The first equality is the chain rule, the second is the gradient definition, and the third is the replicator speed identity. At an interior point, ascent vanishes exactly when all fitness components agree, which is the equilibrium condition.

The scalar potential supplies a monotone quantity for the dynamics. Near a strict local maximum, $\Phi(q)-\Phi(p)$ is a natural Lyapunov candidate. Compact sublevels and isolated critical points permit the stability arguments of Chapter 6 to be applied.

If $M$ is symmetric, the quadratic potential
\begin{equation}\label{eq:symmetricpotential}
\Phi(\p)=\frac12\p^TM\p
\end{equation}
has ambient derivative $M\p$. Thus every symmetric-matrix replicator system is natural-gradient ascent for this potential. Its mean fitness is $\p^TM\p=2\Phi$, so in this class $\bar f'=2\Var_\p(f)$.

\begin{example}[A concave interaction potential]
Let $M=-I$ for three types. Find the dynamics and their velocity at $(0.6,0.3,0.1)$.
\begin{solution}
The potential is $\Phi=-\tfrac12\sum_i p_i^2$, with fitnesses $f_i=-p_i$. Writing $S=\sum_i p_i^2$, the replicator equation is
\[
p_i'=p_i(S-p_i).
\]
At the given state $S=0.46$, so $\p'=(-0.084,0.048,0.036)$. Its ascent rate is $0.0324$. The uniform law uniquely maximizes this strictly concave potential. A KL Lyapunov function can establish convergence of every interior trajectory, as an exercise below shows.
\end{solution}
\end{example}

\section{Stability at a nondegenerate potential maximum}\index{Hessian}
Suppose $q$ is an interior critical point of a smooth potential, expressed in independent coordinates. Write $G(q)$ for its metric matrix and $H(q)$ for the coordinate Hessian of the potential. The gradient field is $F=G^{-1}\nabla\Phi$. Differentiating at a critical point eliminates all terms that multiply $\nabla\Phi(q)=0$, leaving
\[
DF(q)=G(q)^{-1}H(q).
\]
Although this product need not be symmetric in the displayed coordinates, it is similar to a symmetric matrix. The positive-definite square root of $G$ gives
\[
G^{1/2}(G^{-1}H)G^{-1/2}=G^{-1/2}HG^{-1/2}.
\]
The right-hand side is symmetric because $H$ is symmetric. The linearized gradient field therefore has real eigenvalues.

If the Hessian is negative definite, then so is $G^{-1/2}HG^{-1/2}$: its quadratic form on $v$ is the Hessian quadratic form on $G^{-1/2}v$. All eigenvalues are negative, giving a locally asymptotically stable equilibrium. A positive-definite Hessian gives repulsion under ascent. An indefinite nonsingular Hessian produces a saddle with increasing and decreasing linearized directions.

\begin{example}[The local rate at the coexistence maximum]
Apply this calculation to $\Phi(x)=3x-2x^2$ at $q=3/4$.
\begin{solution}
The inverse Bernoulli metric is $q(1-q)=3/16$, and the Hessian is the scalar $-4$. Their product is $-3/4$, reproducing the local rate obtained from the original replicator equation. The potential determines the sign of attraction, while the metric also determines its rate.
\end{solution}
\end{example}

This criterion distinguishes two ways in which trajectories can approach a stable state. A nondegenerate gradient maximum has real linearized decay rates. The cyclic family with $s<1$ has complex linearized eigenvalues and spiraling decay. Its equilibrium attracts, but its rotating linearized motion cannot be the gradient of a smooth potential for a positive-definite metric near that critical point. Convergence by itself is insufficient to identify gradient structure.

The Hessian test is local. It determines behavior near a nondegenerate critical point, while a global convergence statement also needs control of other critical points and of trajectories approaching the boundary. The explicit KL and entropy flows below permit both local and global behavior to be calculated directly.

\section{Linear potentials and exponential evolution}
For $\Phi=\sum_i c_ip_i$, the fitnesses are constants $c_i$. Its natural-gradient flow is the exponential-normalization solution
\[
p_i(t)=\frac{p_i(0)e^{c_it}}{\sum_jp_j(0)e^{c_jt}}.
\]
The formula from Chapter 5 therefore describes an exact optimization flow. The potential is mean constant fitness, and its rate of increase is the variance of the fixed fitness values.

This flow generally approaches the boundary where the highest-fitness types are supported. Positivity at finite time and convergence to a boundary optimum are compatible: a solution can have a positive coordinate at every finite time while that coordinate tends to zero asymptotically.

\section{A KL potential gives mixture contraction}\index{KL divergence!gradient flow}
Fix a positive target $\q$ and choose
\[
\Phi(\p)=-D(\q\|\p).
\]
Its ambient partial derivatives are $f_i=q_i/p_i$, whose weighted mean is 1. The natural-gradient equation simplifies to
\begin{equation}\label{eq:klgradient}
\boxed{p_i'=q_i-p_i.}
\end{equation}
The exact solution is
\begin{equation}\label{eq:klcontraction}
\p(t)=\q+e^{-t}(\p(0)-\q).
\end{equation}
It follows a mixture segment and converges exponentially in probability coordinates. Its geometric route is generally different from the Fisher geodesic connecting the initial law to the target.

The decay identity is
\begin{equation}\label{eq:klgradientdecay}
\frac{d}{dt}D(\q\|\p(t))=-\sum_i\frac{(q_i-p_i)^2}{p_i}.
\end{equation}
For example, direct differentiation gives $1-\sum_iq_i^2/p_i$, and expanding the squared expression proves the equality.

\begin{example}[A target-seeking flow]
Take $\q=(0.3,0.4,0.3)$ and $\p(0)=(0.65,0.2,0.15)$. Find $\p(1)$.
\begin{solution}
The initial displacement is $(0.35,-0.2,-0.15)$. Therefore
\[
\p(1)=(0.3+0.35e^{-1},\ 0.4-0.2e^{-1},\ 0.3-0.15e^{-1}).
\]
Each displacement coordinate contracts by $e^{-1}$, and their sum remains zero.
\end{solution}
\end{example}

The original replicator fitnesses $q_i/p_i$ are defined only in the interior and become singular at a boundary where $q_i>0$. Consequently the smooth-fitness invariant-face theorem from Chapter 3 does not apply to them at that boundary. The simplified vector field $\q-\p$ has a smooth extension that can introduce previously absent types. Specifying the domain and regularity resolves this boundary behavior.

\section{The reverse KL potential and entropy ascent}\index{entropy!gradient flow}
Reverse the directed comparison and let $\Phi(\p)=-D(\p\|\q)$, with the same positive fixed target. Its derivatives are $-\ln(p_i/q_i)-1$, and the flow is
\[
p_i'=p_i\left[-\ln\frac{p_i}{q_i}+D(\p\|\q)\right].
\]
Subtracting the equations for logarithmic ratios gives
\[
\frac{d}{dt}\ln\frac{p_i/q_i}{p_j/q_j}
=-\ln\frac{p_i/q_i}{p_j/q_j}.
\]
Every such log-ratio decays exponentially. For positive initial laws, normalization gives the solution
\begin{equation}\label{eq:reverseklflow}
p_i(t)=\frac{q_i^{1-e^{-t}}p_i(0)^{e^{-t}}}
{\sum_jq_j^{1-e^{-t}}p_j(0)^{e^{-t}}}.
\end{equation}
It converges to $\q$ by a normalized geometric interpolation. The two KL orientations therefore generate different target-seeking paths under the same metric.

For the uniform target $U=(1/n,\ldots,1/n)$, $-D(\p\|U)=H(\p)-\ln n$. Adding a constant to a potential leaves its gradient unchanged, so this is entropy ascent. Directly,
\begin{equation}\label{eq:entropyflow}
p_i'=p_i[-\ln p_i-H(\p)],\qquad
p_i(t)=\frac{p_i(0)^{e^{-t}}}{\sum_jp_j(0)^{e^{-t}}}.
\end{equation}
Entropy increases along this particular dynamical system because it is the generating potential. All initially positive laws approach the uniform distribution.

\begin{example}[Entropy ascent from unequal probabilities]
Let $\p(0)=(1/2,1/3,1/6)$. Find the entropy-ascent state at $t=\ln2$.
\begin{solution}
Here $e^{-t}=1/2$, so the probabilities are proportional to $(\sqrt{1/2},\sqrt{1/3},\sqrt{1/6})$. Normalizing this vector gives the state. At later times the exponent decreases toward zero, and the three positive weights approach equality.
\end{solution}
\end{example}

\section{Cyclic fields and the potential condition}
For the balanced cyclic fitnesses,
\[
f_A=p_B-p_C,\quad f_B=p_C-p_A,\quad f_C=p_A-p_B,
\]
Chapter 3 established nonstationary periodic interior trajectories. If this field were the Fisher gradient of a single-valued potential, then over one period $T$,
\[
0=\Phi(\p(T))-\Phi(\p(0))
=\int_0^T\|\p'(t)\|_F^2\dd t>0,
\]
a contradiction. The Fisher metric still measures its speed and angles; its dynamics circulate without being potential ascent.

Multivariable calculus also supplies a local test. In coordinates $x=p_A,y=p_B$, the differential that a putative potential would need is
\[
d\Phi=(f_A-f_C)dx+(f_B-f_C)dy.
\]
For the balanced cyclic field, this is $(3y-1)dx+(1-3x)dy$. Its cross derivatives are $3$ and $-3$, so it cannot be a differential of a twice differentiable scalar potential on any open region. For general smooth fields on the convex simplex interior, equality of the corresponding cross derivatives is the integrability condition for such a potential.

\begin{figure}[htbp]\centering\begin{tikzpicture}
\begin{scope}[xshift=0.0cm]\begin{scope}[scale=5.7]
\fill[soft] (0,0)--(1,0)--(.5,.8660254)--cycle;
\draw[ink!75] (0,0)--(1,0)--(.5,.8660254)--cycle;
\node[font=\small] at (.5,1.08){Quadratic potential};
\begin{scope}\clip (0,0)--(1,0)--(.5,.8660254)--cycle;
\draw[ink!22] (.5,.2886751) circle(0.08);
\draw[ink!22] (.5,.2886751) circle(0.16);
\draw[ink!22] (.5,.2886751) circle(0.24);
\draw[ink!22] (.5,.2886751) circle(0.32);
\end{scope}
\draw[->,>=Stealth,ink!35,line width=.4pt] (0.5000000,0.6495191)--(0.5000000,0.6175191);
\draw[->,>=Stealth,ink!35,line width=.4pt] (0.5625000,0.5412659)--(0.5647117,0.5093424);
\draw[->,>=Stealth,ink!35,line width=.4pt] (0.6250000,0.4330127)--(0.6161248,0.4022681);
\draw[->,>=Stealth,ink!35,line width=.4pt] (0.6875000,0.3247595)--(0.6564368,0.3324457);
\draw[->,>=Stealth,ink!35,line width=.4pt] (0.7500000,0.2165064)--(0.7234593,0.2343835);
\draw[->,>=Stealth,ink!35,line width=.4pt] (0.8125000,0.1082532)--(0.7847872,0.1242532);
\draw[->,>=Stealth,ink!35,line width=.4pt] (0.4375000,0.5412659)--(0.4352883,0.5093424);
\draw[->,>=Stealth,ink!35,line width=.4pt] (0.5000000,0.4330127)--(0.5000000,0.4010127);
\draw[->,>=Stealth,ink!35,line width=.4pt] (0.5625000,0.3247595)--(0.5347872,0.3087595);
\draw[->,>=Stealth,ink!35,line width=.4pt] (0.6250000,0.2165064)--(0.5972872,0.2325064);
\draw[->,>=Stealth,ink!35,line width=.4pt] (0.6875000,0.1082532)--(0.6587476,0.1222995);
\draw[->,>=Stealth,ink!35,line width=.4pt] (0.3750000,0.4330127)--(0.3838752,0.4022681);
\draw[->,>=Stealth,ink!35,line width=.4pt] (0.4375000,0.3247595)--(0.4652128,0.3087595);
\draw[->,>=Stealth,ink!35,line width=.4pt] (0.5000000,0.2165064)--(0.5000000,0.2485064);
\draw[->,>=Stealth,ink!35,line width=.4pt] (0.5625000,0.1082532)--(0.5403120,0.1313116);
\draw[->,>=Stealth,ink!35,line width=.4pt] (0.3125000,0.3247595)--(0.3435632,0.3324457);
\draw[->,>=Stealth,ink!35,line width=.4pt] (0.3750000,0.2165064)--(0.4027128,0.2325064);
\draw[->,>=Stealth,ink!35,line width=.4pt] (0.4375000,0.1082532)--(0.4596880,0.1313116);
\draw[->,>=Stealth,ink!35,line width=.4pt] (0.2500000,0.2165064)--(0.2765407,0.2343835);
\draw[->,>=Stealth,ink!35,line width=.4pt] (0.3125000,0.1082532)--(0.3412524,0.1222995);
\draw[->,>=Stealth,ink!35,line width=.4pt] (0.1875000,0.1082532)--(0.2152128,0.1242532);
\draw[teal,line width=1.1pt] plot coordinates {(0.2750000,0.1299038) (0.2831183,0.1342287) (0.2910410,0.1384956) (0.2987586,0.1426986) (0.3062639,0.1468326) (0.3135513,0.1508932) (0.3206171,0.1548769) (0.3274591,0.1587805) (0.3340765,0.1626016) (0.3404698,0.1663385) (0.3466406,0.1699898) (0.3525914,0.1735546) (0.3583254,0.1770324) (0.3638468,0.1804230) (0.3691599,0.1837266) (0.3742699,0.1869436) (0.3791821,0.1900748) (0.3839022,0.1931211) (0.3884360,0.1960834) (0.3927894,0.1989630) (0.3969685,0.2017613) (0.4009794,0.2044796) (0.4048281,0.2071195) (0.4085207,0.2096826) (0.4120630,0.2121705) (0.4154609,0.2145848) (0.4187200,0.2169273) (0.4218460,0.2191997) (0.4248442,0.2214036) (0.4277198,0.2235409) (0.4304780,0.2256132) (0.4331237,0.2276223) (0.4356616,0.2295697) (0.4380962,0.2314573) (0.4404320,0.2332866) (0.4426731,0.2350593) (0.4448236,0.2367769) (0.4468875,0.2384411) (0.4488683,0.2400533) (0.4507697,0.2416151) (0.4525951,0.2431280) (0.4543478,0.2445933) (0.4560308,0.2460125) (0.4576472,0.2473869) (0.4591999,0.2487180) (0.4606914,0.2500069) (0.4621245,0.2512550) (0.4635016,0.2524635) (0.4648251,0.2536336) (0.4660973,0.2547665) (0.4673203,0.2558634) (0.4684962,0.2569252) (0.4696270,0.2579532) (0.4707145,0.2589483) (0.4717606,0.2599116) (0.4727669,0.2608440) (0.4737352,0.2617466) (0.4746669,0.2626203) (0.4755636,0.2634659) (0.4764267,0.2642843) (0.4772575,0.2650764) (0.4780574,0.2658431) (0.4788277,0.2665851) (0.4795694,0.2673032) (0.4802838,0.2679981) (0.4809719,0.2686707) (0.4816348,0.2693215) (0.4822734,0.2699513) (0.4828888,0.2705608) (0.4834818,0.2711506) (0.4840533,0.2717213) (0.4846042,0.2722735) (0.4851352,0.2728079) (0.4856471,0.2733250) (0.4861406,0.2738253) (0.4866165,0.2743094) (0.4870754,0.2747778) (0.4875180,0.2752310) (0.4879448,0.2756696) (0.4883565,0.2760938) (0.4887537,0.2765044) (0.4891368,0.2769015) (0.4895065,0.2772858) (0.4898632,0.2776576) (0.4902073,0.2780173) (0.4905394,0.2783653) (0.4908600,0.2787020) (0.4911693,0.2790277) (0.4914679,0.2793429) (0.4917561,0.2796478) (0.4920343,0.2799427) (0.4923029,0.2802281) (0.4925622,0.2805041) (0.4928125,0.2807712) (0.4930542,0.2810295) (0.4932876,0.2812795) (0.4935130,0.2815213) (0.4937307,0.2817552) (0.4939408,0.2819814) (0.4941439,0.2822003) (0.4943399,0.2824121) (0.4945293,0.2826169) (0.4947122,0.2828151) (0.4948889,0.2830068) (0.4950596,0.2831922) (0.4952245,0.2833716) (0.4953838,0.2835451) (0.4955377,0.2837130) (0.4956864,0.2838754) (0.4958300,0.2840325) (0.4959688,0.2841844) (0.4961029,0.2843314) (0.4962325,0.2844735) (0.4963578,0.2846111) (0.4964788,0.2847441) (0.4965957,0.2848728) (0.4967087,0.2849973) (0.4968179,0.2851177) (0.4969235,0.2852342) (0.4970255,0.2853468) (0.4971241,0.2854558) (0.4972194,0.2855612) (0.4973115,0.2856632) (0.4974005,0.2857618) (0.4974866,0.2858572) (0.4975697,0.2859495) (0.4976501,0.2860388) (0.4977278,0.2861251) (0.4978030,0.2862086) (0.4978756,0.2862894) (0.4979458,0.2863676) (0.4980136,0.2864432) (0.4980792,0.2865163) (0.4981427,0.2865870) (0.4982040,0.2866554) (0.4982633,0.2867216) (0.4983206,0.2867856) (0.4983760,0.2868475) (0.4984295,0.2869073) (0.4984813,0.2869653) (0.4985314,0.2870213) (0.4985798,0.2870755) (0.4986266,0.2871279) (0.4986719,0.2871786) (0.4987156,0.2872276) (0.4987579,0.2872750) (0.4987989,0.2873209) (0.4988384,0.2873653) (0.4988767,0.2874082) (0.4989136,0.2874497) (0.4989494,0.2874899) (0.4989840,0.2875287) (0.4990174,0.2875663) (0.4990497,0.2876026) (0.4990810,0.2876378) (0.4991112,0.2876718) (0.4991405,0.2877047) (0.4991687,0.2877365) (0.4991961,0.2877672) (0.4992225,0.2877970) (0.4992481,0.2878258) (0.4992728,0.2878536) (0.4992967,0.2878805) (0.4993198,0.2879066) (0.4993422,0.2879317) (0.4993638,0.2879561) (0.4993847,0.2879797) (0.4994049,0.2880025) (0.4994245,0.2880245) (0.4994434,0.2880458) (0.4994617,0.2880665) (0.4994794,0.2880864) (0.4994965,0.2881057) (0.4995130,0.2881244) (0.4995290,0.2881424) (0.4995445,0.2881599) (0.4995594,0.2881768) (0.4995739,0.2881931) (0.4995879,0.2882089) (0.4996014,0.2882242) (0.4996145,0.2882390) (0.4996272,0.2882533) (0.4996394,0.2882671) (0.4996512,0.2882805) (0.4996627,0.2882934) (0.4996738,0.2883059) (0.4996845,0.2883180) (0.4996948,0.2883297) (0.4997048,0.2883411) (0.4997145,0.2883520) (0.4997239,0.2883626) (0.4997330,0.2883728) (0.4997417,0.2883828) (0.4997502,0.2883923) (0.4997584,0.2884016) (0.4997663,0.2884106) (0.4997740,0.2884192) (0.4997814,0.2884276) (0.4997886,0.2884357) (0.4997955,0.2884436) (0.4998022,0.2884512) (0.4998087,0.2884585) (0.4998150,0.2884656) (0.4998211,0.2884725) (0.4998269,0.2884791) (0.4998326,0.2884856) (0.4998381,0.2884918) (0.4998434,0.2884978) (0.4998485,0.2885036) (0.4998535,0.2885092) (0.4998583,0.2885147) (0.4998630,0.2885199) (0.4998675,0.2885250) (0.4998718,0.2885299) (0.4998760,0.2885347) (0.4998801,0.2885393) (0.4998840,0.2885438) (0.4998878,0.2885481) (0.4998915,0.2885522) (0.4998951,0.2885563) (0.4998985,0.2885602) (0.4999018,0.2885639) (0.4999050,0.2885676) (0.4999082,0.2885711) (0.4999112,0.2885745) (0.4999141,0.2885778) (0.4999169,0.2885810) (0.4999196,0.2885841) (0.4999223,0.2885871) (0.4999248,0.2885900) (0.4999273,0.2885927) (0.4999297,0.2885954) (0.4999320,0.2885981) (0.4999342,0.2886006) (0.4999364,0.2886030) (0.4999384,0.2886054) (0.4999405,0.2886077) (0.4999424,0.2886099) (0.4999443,0.2886120) (0.4999461,0.2886141) (0.4999479,0.2886161) (0.4999496,0.2886180) (0.4999513,0.2886199) (0.4999529,0.2886217) (0.4999544,0.2886235) (0.4999559,0.2886252) (0.4999573,0.2886268) (0.4999587,0.2886284) (0.4999601,0.2886299) (0.4999614,0.2886314) (0.4999627,0.2886328) (0.4999639,0.2886342) (0.4999651,0.2886356) (0.4999662,0.2886369) (0.4999673,0.2886381) (0.4999684,0.2886393) (0.4999694,0.2886405) (0.4999704,0.2886416) (0.4999714,0.2886427) (0.4999723,0.2886438) (0.4999733,0.2886448) (0.4999741,0.2886458) (0.4999750,0.2886468) (0.4999758,0.2886477) (0.4999766,0.2886486) (0.4999774,0.2886495) (0.4999781,0.2886503) (0.4999788,0.2886511) (0.4999795,0.2886519) (0.4999802,0.2886527) (0.4999808,0.2886534) (0.4999815,0.2886541) (0.4999821,0.2886548) (0.4999827,0.2886555) (0.4999832,0.2886561) (0.4999838,0.2886568) (0.4999843,0.2886574) (0.4999848,0.2886579) (0.4999853,0.2886585) (0.4999858,0.2886590) (0.4999863,0.2886596) (0.4999867,0.2886601) (0.4999872,0.2886606) (0.4999876,0.2886611) (0.4999880,0.2886615) (0.4999884,0.2886620) (0.4999888,0.2886624) (0.4999891,0.2886628) (0.4999895,0.2886632) (0.4999898,0.2886636) (0.4999902,0.2886640) (0.4999905,0.2886643) (0.4999908,0.2886647) (0.4999911,0.2886650) (0.4999914,0.2886654) (0.4999917,0.2886657) (0.4999919,0.2886660) (0.4999922,0.2886663) (0.4999925,0.2886666) (0.4999927,0.2886669) (0.4999929,0.2886671) (0.4999932,0.2886674) (0.4999934,0.2886677) (0.4999936,0.2886679) (0.4999938,0.2886681) (0.4999940,0.2886684) (0.4999942,0.2886686) (0.4999944,0.2886688) (0.4999946,0.2886690) (0.4999948,0.2886692) (0.4999949,0.2886694) (0.4999951,0.2886696) (0.4999953,0.2886698) (0.4999954,0.2886700) (0.4999956,0.2886701) (0.4999957,0.2886703) (0.4999959,0.2886704) (0.4999960,0.2886706) (0.4999961,0.2886707) (0.4999963,0.2886709) (0.4999964,0.2886710) (0.4999965,0.2886712) (0.4999966,0.2886713) (0.4999967,0.2886714) (0.4999968,0.2886715) (0.4999969,0.2886717) (0.4999970,0.2886718) (0.4999971,0.2886719) (0.4999972,0.2886720) (0.4999973,0.2886721) (0.4999974,0.2886722) (0.4999975,0.2886723) (0.4999976,0.2886724) (0.4999977,0.2886725) (0.4999977,0.2886726) (0.4999978,0.2886726) (0.4999979,0.2886727) (0.4999979,0.2886728) (0.4999980,0.2886729) (0.4999981,0.2886730) (0.4999981,0.2886730) (0.4999982,0.2886731) (0.4999983,0.2886732) (0.4999983,0.2886732) (0.4999984,0.2886733) (0.4999984,0.2886734) (0.4999985,0.2886734) (0.4999985,0.2886735) (0.4999986,0.2886735) (0.4999986,0.2886736) (0.4999987,0.2886736)};
\fill[copper] (0.2750000,0.1299038) circle(.009);
\end{scope}\end{scope}
\begin{scope}[xshift=7.4cm]\begin{scope}[scale=5.7]
\fill[soft] (0,0)--(1,0)--(.5,.8660254)--cycle;
\draw[ink!75] (0,0)--(1,0)--(.5,.8660254)--cycle;
\node[font=\small] at (.5,1.08){Balanced cyclic model};
\draw[->,>=Stealth,ink!35,line width=.4pt] (0.5000000,0.6495191)--(0.5320000,0.6495191);
\draw[->,>=Stealth,ink!35,line width=.4pt] (0.5625000,0.5412659)--(0.5876429,0.5214710);
\draw[->,>=Stealth,ink!35,line width=.4pt] (0.6250000,0.4330127)--(0.6437266,0.4070644);
\draw[->,>=Stealth,ink!35,line width=.4pt] (0.6875000,0.3247595)--(0.7006086,0.2955677);
\draw[->,>=Stealth,ink!35,line width=.4pt] (0.7500000,0.2165064)--(0.7545714,0.1848346);
\draw[->,>=Stealth,ink!35,line width=.4pt] (0.8125000,0.1082532)--(0.7965000,0.0805404);
\draw[->,>=Stealth,ink!35,line width=.4pt] (0.4375000,0.5412659)--(0.4626429,0.5610607);
\draw[->,>=Stealth,ink!35,line width=.4pt] (0.5000000,0.4330127)--(0.5320000,0.4330127);
\draw[->,>=Stealth,ink!35,line width=.4pt] (0.5625000,0.3247595)--(0.5785000,0.2970467);
\draw[->,>=Stealth,ink!35,line width=.4pt] (0.6250000,0.2165064)--(0.6090000,0.1887935);
\draw[->,>=Stealth,ink!35,line width=.4pt] (0.6875000,0.1082532)--(0.6577857,0.0963763);
\draw[->,>=Stealth,ink!35,line width=.4pt] (0.3750000,0.4330127)--(0.3937266,0.4589610);
\draw[->,>=Stealth,ink!35,line width=.4pt] (0.4375000,0.3247595)--(0.4535000,0.3524723);
\draw[->,>=Stealth,ink!35,line width=.4pt] (0.5000000,0.2165064)--(0.4680000,0.2165064);
\draw[->,>=Stealth,ink!35,line width=.4pt] (0.5625000,0.1082532)--(0.5306648,0.1050096);
\draw[->,>=Stealth,ink!35,line width=.4pt] (0.3125000,0.3247595)--(0.3256086,0.3539514);
\draw[->,>=Stealth,ink!35,line width=.4pt] (0.3750000,0.2165064)--(0.3590000,0.2442192);
\draw[->,>=Stealth,ink!35,line width=.4pt] (0.4375000,0.1082532)--(0.4056648,0.1114967);
\draw[->,>=Stealth,ink!35,line width=.4pt] (0.2500000,0.2165064)--(0.2545714,0.2481781);
\draw[->,>=Stealth,ink!35,line width=.4pt] (0.3125000,0.1082532)--(0.2827857,0.1201301);
\draw[->,>=Stealth,ink!35,line width=.4pt] (0.1875000,0.1082532)--(0.1715000,0.1359660);
\draw[teal,line width=1.1pt] plot coordinates {(0.2750000,0.1299038) (0.2687861,0.1359693) (0.2633953,0.1424833) (0.2588263,0.1494582) (0.2550746,0.1569049) (0.2521329,0.1648330) (0.2499910,0.1732496) (0.2486364,0.1821597) (0.2480539,0.1915650) (0.2482262,0.2014642) (0.2491336,0.2118520) (0.2507542,0.2227190) (0.2530639,0.2340513) (0.2560364,0.2458302) (0.2596432,0.2580320) (0.2638539,0.2706276) (0.2686360,0.2835827) (0.2739548,0.2968578) (0.2797744,0.3104083) (0.2860570,0.3241845) (0.2927636,0.3381324) (0.2998546,0.3521938) (0.3072893,0.3663073) (0.3150274,0.3804085) (0.3230284,0.3944310) (0.3312529,0.4083070) (0.3396626,0.4219685) (0.3482206,0.4353475) (0.3568925,0.4483772) (0.3656462,0.4609926) (0.3744523,0.4731313) (0.3832848,0.4847336) (0.3921209,0.4957435) (0.4009413,0.5061090) (0.4097304,0.5157821) (0.4184761,0.5247190) (0.4271700,0.5328806) (0.4358070,0.5402320) (0.4443852,0.5467428) (0.4529057,0.5523867) (0.4613725,0.5571418) (0.4697916,0.5609897) (0.4781711,0.5639162) (0.4865209,0.5659106) (0.4948517,0.5669656) (0.5031752,0.5670775) (0.5115033,0.5662458) (0.5198475,0.5644736) (0.5282191,0.5617671) (0.5366280,0.5581364) (0.5450828,0.5535949) (0.5535902,0.5481598) (0.5621545,0.5418522) (0.5707776,0.5346974) (0.5794583,0.5267249) (0.5881924,0.5179683) (0.5969722,0.5084660) (0.6057865,0.4982605) (0.6146204,0.4873987) (0.6234557,0.4759319) (0.6322703,0.4639151) (0.6410390,0.4514071) (0.6497331,0.4384698) (0.6583211,0.4251676) (0.6667692,0.4115670) (0.6750410,0.3977356) (0.6830985,0.3837417) (0.6909025,0.3696531) (0.6984128,0.3555366) (0.7055889,0.3414571) (0.7123903,0.3274769) (0.7187770,0.3136550) (0.7247097,0.3000464) (0.7301506,0.2867018) (0.7350630,0.2736670) (0.7394122,0.2609829) (0.7431650,0.2486850) (0.7462905,0.2368035) (0.7487598,0.2253635) (0.7505460,0.2143848) (0.7516244,0.2038824) (0.7519728,0.1938668) (0.7515709,0.1843441) (0.7504007,0.1753167) (0.7484467,0.1667834) (0.7456955,0.1587402) (0.7421366,0.1511801) (0.7377618,0.1440944) (0.7325660,0.1374722) (0.7265472,0.1313016) (0.7197065,0.1255693) (0.7120491,0.1202615) (0.7035838,0.1153640) (0.6943240,0.1108622) (0.6842874,0.1067420) (0.6734970,0.1029891) (0.6619809,0.0995902) (0.6497724,0.0965323) (0.6369106,0.0938032) (0.6234402,0.0913916) (0.6094113,0.0892873) (0.5948793,0.0874811) (0.5799050,0.0859648) (0.5645535,0.0847314) (0.5488941,0.0837752) (0.5329997,0.0830916) (0.5169458,0.0826775) (0.5008096,0.0825308) (0.4846695,0.0826507) (0.4686038,0.0830380) (0.4526900,0.0836944) (0.4370037,0.0846231) (0.4216182,0.0858284) (0.4066030,0.0873161) (0.3920239,0.0890930) (0.3779419,0.0911669) (0.3644132,0.0935472) (0.3514883,0.0962440) (0.3392125,0.0992685) (0.3276253,0.1026326) (0.3167608,0.1063494) (0.3066476,0.1104322) (0.2973089,0.1148950) (0.2887633,0.1197523) (0.2810244,0.1250183) (0.2741016,0.1307074) (0.2680003,0.1368334) (0.2627220,0.1434097) (0.2582651,0.1504485) (0.2546247,0.1579606) (0.2517930,0.1659551) (0.2497596,0.1744390) (0.2485115,0.1834167) (0.2480333,0.1928898) (0.2483073,0.2028562) (0.2493136,0.2133102) (0.2510299,0.2242417) (0.2534319,0.2356363) (0.2564929,0.2474746) (0.2601842,0.2597320) (0.2644750,0.2723789) (0.2693325,0.2853803) (0.2747220,0.2986958) (0.2806071,0.3122802) (0.2869500,0.3260832) (0.2937116,0.3400501) (0.3008520,0.3541224) (0.3083306,0.3682380) (0.3161070,0.3823324) (0.3241409,0.3963389) (0.3323930,0.4101896) (0.3408250,0.4238164) (0.3494006,0.4371516) (0.3580856,0.4501284) (0.3668480,0.4626823) (0.3756593,0.4747510) (0.3844935,0.4862755) (0.3933285,0.4972003) (0.4021455,0.5074738) (0.4109293,0.5170487) (0.4196683,0.5258818) (0.4283547,0.5339345) (0.4369836,0.5411727) (0.4455538,0.5475664) (0.4540667,0.5530900) (0.4625266,0.5577219) (0.4709397,0.5614446) (0.4793146,0.5642442) (0.4876611,0.5661105) (0.4959902,0.5670366) (0.5043137,0.5670194) (0.5126432,0.5660588) (0.5209906,0.5641584) (0.5293667,0.5613249) (0.5377814,0.5575687) (0.5462431,0.5529039) (0.5547580,0.5473481) (0.5633304,0.5409231) (0.5719615,0.5336546) (0.5806498,0.5255728) (0.5893907,0.5167119)};
\fill[copper] (0.2750000,0.1299038) circle(.009);
\end{scope}\end{scope}
\end{tikzpicture}\caption{Natural-gradient ascent for $\Phi=-\tfrac12\sum_i p_i^2$ and balanced cyclic dynamics from the same initial law. Left contours are potential levels. Short arrows show field directions with normalized drawing lengths.}\end{figure}

For numerical motion under a smooth fitness field, a useful finite update is
\[
p_i^{\mathrm{new}}=\frac{p_i e^{\eta f_i(\p)}}{\sum_jp_j e^{\eta f_j(\p)}}.
\]
It stays positive and normalized for finite scores. Expanding in small $\eta$ yields $p_i^{\mathrm{new}}=p_i+\eta p_i(f_i-\bar f)+O(\eta^2)$, so it approximates a replicator step. It is exact for a constant-fitness time step. For a varying fitness field, preservation of the simplex alone does not establish a particular numerical error or monotonicity guarantee.

\FloatBarrier
\section{Exercises}
\begin{exercises}
\item Derive the coordinate gradient formula $G^{-1}\nabla\Phi$ from the defining inner-product identity.
\item Use Cauchy--Schwarz to prove that the normalized metric gradient maximizes directional increase among unit tangent vectors.
\item For $\Phi(x)=2x-(5/2)x^2$, derive the Bernoulli natural-gradient flow and classify its equilibria.
\item Verify \eqref{eq:simplexnatural}, including the tangent constraint and independence from the chosen ambient extension.
\item At $\p=(1/2,1/3,1/6)$ and $\Phi=2p_A+p_B$, compute the natural gradient and its ascent rate.
\item For a symmetric payoff matrix, prove that $\tfrac12\p^TM\p$ is a potential and that mean fitness increases at rate $2\Var_\p(f)$.
\item For $M=-I$ with $n$ types, use the multi-strategy KL derivative from Chapter 6 and a uniform reference to prove global interior convergence to the uniform law.
\item Derive the solution \eqref{eq:klcontraction} and verify the exact divergence decay identity \eqref{eq:klgradientdecay}.
\item Compute the target-seeking state in the worked example at $t=1$ numerically. Determine the time needed to reduce every coordinate displacement by a factor of 10.
\item Explain why the smooth extension $p'=q-p$ can introduce a type that was absent initially without contradicting the smooth replicator invariant-face theorem.
\item Derive the reverse-KL gradient and the normalized geometric solution \eqref{eq:reverseklflow}.
\item Derive the entropy gradient and show that its log-probability ratios decay exponentially.
\item For entropy ascent from $(1/2,1/3,1/6)$, find the normalized state at $t=\ln2$ to four decimal places.
\item Prove that no nonstationary periodic trajectory can be a gradient-ascent trajectory for any positive-definite Riemannian metric and a single-valued potential.
\item In reduced simplex coordinates, derive the differential coefficients $f_A-f_C$ and $f_B-f_C$. Apply the cross-derivative test to balanced cyclic competition.
\item Expand the exponential finite-step update through first order in $\eta$ and recover the replicator field.
\item \emph{Exploration.} Compare a Fisher geodesic to the mixture-contraction trajectory with the same initial and target laws. Measure route length separately from the rate of potential improvement.
\item \emph{Exploration.} Compare linear-potential ascent, entropy ascent, and balanced cyclic motion from the same positive initial state. Track a justified potential or conserved quantity for each system.
\end{exercises}

\clearpage\backmatter
\chapter*{Selected Solutions}
\section*{Chapter 9}
\begin{answers}
\item[9.3] The flow is $x'=x(1-x)(2-5x)$. The equilibria are $0,2/5,1$, with the interior state attracting and the endpoints repelling from the interior.
\item[9.5] The derivative fitnesses are $(2,1,0)$ and their mean is $4/3$. The gradient is $(1/3,-1/9,-2/9)$. Its squared Fisher norm, and ascent rate, is $5/9$.
\item[9.7] With $S=\sum_i p_i^2$, the KL derivative is $1/n-S\leq0$, strictly below zero away from uniform. The divergence from a positive uniform reference tends to infinity at the boundary, so its sublevels are compact in the interior. The strict Lyapunov argument gives convergence.
\item[9.9] The state is approximately $(0.4288,0.3264,0.2448)$. Since the displacement factor is $e^{-t}$, a tenfold reduction takes $t=\ln10$.
\item[9.11] The ambient derivatives are $-\ln(p_i/q_i)-1$. Subtract their weighted mean to obtain the displayed vector field. Log-ratios of $p_i/q_i$ obey $r'=-r$; integrating and normalizing gives \eqref{eq:reverseklflow}.
\item[9.13] Normalize $(\sqrt{1/2},\sqrt{1/3},\sqrt{1/6})$. The resulting probabilities are approximately
\[
(0.4177,0.3411,0.2412).
\]
\item[9.15] A tangent vector has $v_C=-v_A-v_B$, so $\sum_i f_iv_i=(f_A-f_C)v_A+(f_B-f_C)v_B$. For the cyclic field the two cross derivatives are $3$ and $-3$, which disagree.
\item[9.16] Expand the numerator as $p_i(1+\eta f_i+O(\eta^2))$ and the denominator as $1+\eta\bar f+O(\eta^2)$. Multiplying by the reciprocal denominator gives $p_i+\eta p_i(f_i-\bar f)+O(\eta^2)$.
\end{answers}
\end{document}
